\documentclass{article}
\usepackage{amsmath,amssymb,amsthm}
\IfFileExists{hyperref.sty}{%
  \IfFileExists{color.sty}{\usepackage{color}\definecolor{paperlinkblue}{rgb}{0.10,0.25,0.42}%
    \usepackage[colorlinks=true,linkcolor=paperlinkblue]{hyperref}}{\usepackage[hidelinks]{hyperref}}%
}{}
\ifdefined\hyperlink
  \newcommand{\paperanchor}[1]{\hypertarget{#1}{}}%
  \newcommand{\paperlink}[2]{\hyperlink{#1}{\begin{NoHyper}#2\end{NoHyper}}}%
\else
  \ifdefined\pdfextension
    \ifnum\outputmode>0
      \def\paperdest{\pdfextension dest}\def\paperstart{\pdfextension startlink}%
      \def\paperend{\pdfextension endlink}\def\papercolor{\pdfextension colorstack}%
    \fi
  \else\ifdefined\pdfstartlink
    \ifnum\pdfoutput>0
      \let\paperdest\pdfdest\let\paperstart\pdfstartlink\let\paperend\pdfendlink\let\papercolor\pdfcolorstack
    \fi
  \fi\fi
  \ifdefined\paperdest
    \newcommand{\paperanchor}[1]{\paperdest name{#1} xyz}%
    \newcommand{\paperlink}[2]{\leavevmode\paperstart attr{/Border [0 0 0]} goto name{#1}%
      \papercolor 0 push {0.10 0.25 0.42 rg 0.10 0.25 0.42 RG}#2\papercolor 0 pop\paperend}%
  \else\ifdefined\XeTeXversion
    \newcommand{\paperanchor}[1]{\special{pdf:dest (#1) [@thispage /XYZ null @ypos null]}}%
    \newcommand{\paperlink}[2]{\leavevmode\special{pdf:bann << /Subtype /Link /Border [0 0 0] /A << /S /GoTo /D (#1) >> >>}%
      \special{color push rgb 0.10 0.25 0.42}#2\special{color pop}\special{pdf:eann}}%
  \else
    \PackageError{lean-to-latex}{No PDF hyperlink backend available}{Use XeLaTeX, pdfLaTeX, LuaLaTeX, or install hyperref.}%
  \fi\fi
\fi
\newcommand{\paperlineanchor}[1]{\raisebox{10pt}[0pt][0pt]{\paperanchor{#1}}}
\DeclareRobustCommand{\paperresultname}[1]{\paperlineanchor{theorem:\arabic{theorem}}#1}
\newtheoremstyle{paperresult}{8pt}{5pt}{\normalfont\raggedright}{}{\bfseries}{.}{\newline}{}
\theoremstyle{paperresult}
\newtheorem{theorem}{\paperresultname{Theorem}}
\newcounter{paperpremise}
\newtheorem{lemma}[theorem]{\paperresultname{Lemma}}
\newtheorem{corollary}[theorem]{\paperresultname{Corollary}}
\newcounter{definingequation}
\newenvironment{proofstep}{\par\smallskip\ifdim\leftskip<3em\relax\advance\leftskip by 1em\relax\fi\noindent\ignorespaces}{\par\smallskip}
\setlength{\textwidth}{6.5in}
\setlength{\oddsidemargin}{0in}
\AddToHook{env/proof/begin}{\raggedright}
\setlength{\emergencystretch}{\linewidth}
\title{Parity of matchings}
\author{}
\date{}
\begin{document}
\maketitle
\section*{Definitions}
\paragraph{Is matching.} \par\noindent\refstepcounter{definingequation}\label{defeq:1}\paperlineanchor{defeq:1}\textup{(\paperlink{defeq:1}{D1})} Let \(V\) be a type. Let \(G\) be a simple graph on \(V\). Let \(M\) be a subgraph of \(G\). Then \((M\text{ is a matching}) = (\forall v\, :\, V,\; v \in V(M) \implies (\exists! w\, :\, V,\; v \sim_{M} w))\).

\subsection*{Prerequisites}
\begin{description}
\item[\paperlineanchor{premise:46696e7365742e73756d5f636f6e6772}\paperlink{premise:46696e7365742e73756d5f636f6e6772}{P1.}] \refstepcounter{paperpremise}\label{premise:46696e7365742e73756d5f636f6e6772}\textbf{Equality of sums with equal summands.} Let \(\iota\) and \(M\) be types. Let \(s_{1}\) and \(s_{2}\) be finite subsets of \(\iota\). Equip \(M\) with a commutative additive monoid structure. Let \(f\) and \(g\) be functions of type \(\iota \to M\). Assume \(s_{1} = s_{2}\). Assume \(\forall x\, :\, \iota,\; x \in s_{2} \implies f\left(x\right) = g\left(x\right)\). Then \(\sum_{i\in s_{1}}\left(f(i)\right) = \sum_{i\in s_{2}}\left(g(i)\right)\).
\item[\paperlineanchor{premise:46696e7365742e73756d5f636f6e7374}\paperlink{premise:46696e7365742e73756d5f636f6e7374}{P2.}] \refstepcounter{paperpremise}\label{premise:46696e7365742e73756d5f636f6e7374}\textbf{The sum of a constant.} Let \(\iota\) and \(M\) be types. Let \(s\) be a finite subset of \(\iota\). Equip \(M\) with a commutative additive monoid structure. Let \(b\) be an element of \(M\). Then \(\sum_{x\in s}\left(b\right) = \sum_{k=1}^{\left|s\right|}\left(b\right)\).
\item[\paperlineanchor{premise:46696e747970652e696e73744661737453756273696e676c65746f6e}\paperlink{premise:46696e747970652e696e73744661737453756273696e676c65746f6e}{P3.}] \refstepcounter{paperpremise}\label{premise:46696e747970652e696e73744661737453756273696e676c65746f6e}\textbf{Inst fast subsingleton.} Let \(\alpha\) be a type. Then \(\operatorname{Fast\,Subsingleton}\left(\operatorname{Fintype}\left(\alpha\right)\right)\).
\item[\paperlineanchor{premise:5365742e746f46696e7365745f63617264}\paperlink{premise:5365742e746f46696e7365745f63617264}{P4.}] \refstepcounter{paperpremise}\label{premise:5365742e746f46696e7365745f63617264}\textbf{To finset cardinality.} Let \(\alpha\) be a type. Let \(s\) be a subset of \(\alpha\). Fix a finite enumeration of \(s\). Then \(\left|\operatorname{as\,finite\,set}(s)\right| = \left|s\right|\).
\item[\paperlineanchor{premise:53696d706c6547726170682e53756267726170682e636f655f646567726565}\paperlink{premise:53696d706c6547726170682e53756267726170682e636f655f646567726565}{P5.}] \refstepcounter{paperpremise}\label{premise:53696d706c6547726170682e53756267726170682e636f655f646567726565}\textbf{Coercion degree.} Let \(V\) be a type. Let \(G\) be a simple graph on \(V\). Let \(G^{'}\) be a subgraph of \(G\). Let \(v\) be an element of \(V(G^{'})\). Fix a finite enumeration of \(N_{\operatorname{graph}(G^{'})}(v)\). Fix a finite enumeration of \(N_{G^{'}}(v)\). Then \(\deg_{\operatorname{graph}(G^{'})}(v) = \deg_{G^{'}}(v)\).
\item[\paperlineanchor{premise:53696d706c6547726170682e53756267726170682e6465677265655f65715f6f6e655f6966665f756e697175655f61646a}\paperlink{premise:53696d706c6547726170682e53756267726170682e6465677265655f65715f6f6e655f6966665f756e697175655f61646a}{P6.}] \refstepcounter{paperpremise}\label{premise:53696d706c6547726170682e53756267726170682e6465677265655f65715f6f6e655f6966665f756e697175655f61646a}\textbf{Degree equals one if and only if unique adj.} Let \(V\) be a type. Let \(G\) be a simple graph on \(V\). Let \(G^{'}\) be a subgraph of \(G\). Let \(v\) be an element of \(V\). Fix a finite enumeration of \(N_{G^{'}}(v)\). Then \(\deg_{G^{'}}(v) = 1 \iff \left(\exists! w\, :\, V,\; v \sim_{G^{'}} w\right)\).
\item[\paperlineanchor{premise:53696d706c6547726170682e73756d5f646567726565735f65715f74776963655f636172645f6564676573}\paperlink{premise:53696d706c6547726170682e73756d5f646567726565735f65715f74776963655f636172645f6564676573}{P7.}] \refstepcounter{paperpremise}\label{premise:53696d706c6547726170682e73756d5f646567726565735f65715f74776963655f636172645f6564676573}\textbf{The degree-sum formula.} Let \(V\) be a type. Let \(G\) be a simple graph on \(V\). Fix a finite enumeration of \(V\). Assume adjacency in \(G\) is decidable. Then \(\sum_{v\in V}\left(\deg_{G}(v)\right) = 2 \cdot \left|E(G)\right|\).
\item[\paperlineanchor{premise:537562747970652e636f655f70726f70}\paperlink{premise:537562747970652e636f655f70726f70}{P8.}] \refstepcounter{paperpremise}\label{premise:537562747970652e636f655f70726f70}\textbf{Coercion prop.} Let \(\alpha\) be a type. Let \(S\) be a subset of \(\alpha\). Let \(a\) be an element of \(S\). Then \(a \in S\).
\item[\paperlineanchor{premise:65715f6f665f686571}\paperlink{premise:65715f6f665f686571}{P9.}] \refstepcounter{paperpremise}\label{premise:65715f6f665f686571}\textbf{Equals of heq.} Let \(\alpha\) be a sort. Let \(a\) and \(a^{'}\) be elements of \(\alpha\). Assume \(\left(a\, :\, \alpha\right)\text{ is heterogeneously equal to }\left(a^{'}\, :\, \alpha\right)\). Then \(a = a^{'}\).
\item[\paperlineanchor{premise:666f72616c6c5f636f6e6772}\paperlink{premise:666f72616c6c5f636f6e6772}{P10.}] \refstepcounter{paperpremise}\label{premise:666f72616c6c5f636f6e6772}\textbf{Forall congruence.} Let \(\alpha\) be a sort. Let \(p\) and \(q\) be functions of type \(\alpha \to \mathsf{Prop}\). Assume \(\forall a\, :\, \alpha,\; p\left(a\right) = q\left(a\right)\). Then \((\forall a\, :\, \alpha,\; p(a)) = (\forall a\, :\, \alpha,\; q(a))\).
\item[\paperlineanchor{premise:6966665f73656c66}\paperlink{premise:6966665f73656c66}{P11.}] \refstepcounter{paperpremise}\label{premise:6966665f73656c66}\textbf{If and only if self.} Let \(p\) be a proposition. Then \(\left(p \iff p\right) = \mathrm{True}\).
\item[\paperlineanchor{premise:696d706c6965735f636f6e6772}\paperlink{premise:696d706c6965735f636f6e6772}{P12.}] \refstepcounter{paperpremise}\label{premise:696d706c6965735f636f6e6772}\textbf{Implies congruence.} Let \(p_{1}\) and \(p_{2}\) be sorts. Let \(q_{1}\) and \(q_{2}\) be sorts. Assume \(p_{1} = p_{2}\). Assume \(q_{1} = q_{2}\). Then \(\left(p_{1} \to q_{1}\right) = \left(p_{2} \to q_{2}\right)\).
\item[\paperlineanchor{premise:6d756c5f6f6e65}\paperlink{premise:6d756c5f6f6e65}{P13.}] \refstepcounter{paperpremise}\label{premise:6d756c5f6f6e65}\textbf{Product one.} Let \(M\) be a type. Fix a structure \(\mathit{inst}\) of type \(\operatorname{Mul\,One\,Class}\left(M\right)\). Let \(a\) be an element of \(M\). Then \(a \cdot 1 = a\).
\item[\paperlineanchor{premise:6f665f65715f74727565}\paperlink{premise:6f665f65715f74727565}{P14.}] \refstepcounter{paperpremise}\label{premise:6f665f65715f74727565}\textbf{Of equals true.} Let \(p\) be a proposition. Assume \(p = \mathrm{True}\). Then \(p\).
\item[\paperlineanchor{premise:74776f5f6d756c}\paperlink{premise:74776f5f6d756c}{P15.}] \refstepcounter{paperpremise}\label{premise:74776f5f6d756c}\textbf{Two product.} Let \(\alpha\) be a type. Fix a structure \(\mathit{inst}\) of type \(\operatorname{Non\,Assoc\,Semiring}\left(\alpha\right)\). Let \(n\) be an element of \(\alpha\). Then \(2 \cdot n = n + n\).
\end{description}

\section{Results}
\begin{lemma}[{The degree-one characterization of a matching}]\label{theorem:1}
Let \(V\) be a type. Let \(G\) be a simple graph on \(V\). Let \(M\) be a subgraph of \(G\). For each \(v \in V\), fix a finite enumeration of its neighbors in \(M\). Then
\begin{center}
\(\displaystyle M\text{ is a matching} \iff (\forall v\, :\, V,\; v \in V(M) \implies \deg_{M}(v) = 1).\)
\end{center}
\end{lemma}

\begin{proof}
For the required implication, argue as follows.
\begin{proofstep}
Fix \(v\) in \(V\).
By definition, we obtain \(\left(v \in V(M)\right) = \left(v \in V(M)\right)\).
By \paperlink{premise:53696d706c6547726170682e53756267726170682e6465677265655f65715f6f6e655f6966665f756e697175655f61646a}{Prerequisite~P\ref{premise:53696d706c6547726170682e53756267726170682e6465677265655f65715f6f6e655f6966665f756e697175655f61646a}}, we obtain \(\deg_{M}(v) = 1 \iff \left(\exists! w\, :\, V,\; v \sim_{M} w\right)\).
By propositional extensionality, we obtain \(\left(\deg_{M}(v) = 1\right) = \left(\exists! w\, :\, V,\; v \sim_{M} w\right)\).
By \paperlink{premise:696d706c6965735f636f6e6772}{Prerequisite~P\ref{premise:696d706c6965735f636f6e6772}} with \(p_{1} = v \in V(M)\) and \(p_{2} = v \in V(M)\) and \(q_{1} = \deg_{M}(v) = 1\) and \(q_{2} = \exists! w\, :\, V,\; v \sim_{M} w\), we obtain \((v \in V(M) \implies \deg_{M}(v) = 1) = (v \in V(M) \implies (\exists! w\, :\, V,\; v \sim_{M} w))\).
\end{proofstep}
By \paperlink{premise:666f72616c6c5f636f6e6772}{Prerequisite~P\ref{premise:666f72616c6c5f636f6e6772}}, we obtain \((\forall a\, :\, V,\; a \in V(M) \implies \deg_{M}(a) = 1) = (\forall a\, :\, V,\; a \in V(M) \implies (\exists! w\, :\, V,\; a \sim_{M} w))\).
Applying the same function to equal arguments, we obtain \(((\forall v\, :\, V,\; v \in V(M) \implies (\exists! w\, :\, V,\; v \sim_{M} w)) \iff (\forall a\, :\, V,\; a \in V(M) \implies \deg_{M}(a) = 1)) = ((\forall v\, :\, V,\; v \in V(M) \implies (\exists! w\, :\, V,\; v \sim_{M} w)) \iff (\forall a\, :\, V,\; a \in V(M) \implies (\exists! w\, :\, V,\; a \sim_{M} w)))\).
By \paperlink{premise:6966665f73656c66}{Prerequisite~P\ref{premise:6966665f73656c66}} with \(p = \forall v\, :\, V,\; v \in V(M) \implies \left(\exists! w\, :\, V,\; v \sim_{M} w\right)\), we obtain \(((\forall v\, :\, V,\; v \in V(M) \implies (\exists! w\, :\, V,\; v \sim_{M} w)) \iff (\forall v\, :\, V,\; v \in V(M) \implies (\exists! w\, :\, V,\; v \sim_{M} w))) = \mathrm{True}\).
Chaining the equalities, we obtain \(((\forall v\, :\, V,\; v \in V(M) \implies (\exists! w\, :\, V,\; v \sim_{M} w)) \iff (\forall a\, :\, V,\; a \in V(M) \implies \deg_{M}(a) = 1)) = \mathrm{True}\).
By \paperlink{premise:6f665f65715f74727565}{Prerequisite~P\ref{premise:6f665f65715f74727565}}, we obtain \(M\text{ is a matching} \iff \left(\forall v\, :\, V,\; v \in V(M) \implies \deg_{M}(v) = 1\right)\).
\end{proof}

\begin{theorem}[{A matching covers an even number of vertices}]\label{theorem:2}
Let \(V\) be a type. Let \(G\) be a simple graph on \(V\). Let \(M\) be a subgraph of \(G\). Fix a finite enumeration of \(V(M)\). Assume \(M\text{ is a matching}\). Then \(\operatorname{Even}\left(\left|\operatorname{as\,finite\,set}(V(M))\right|\right)\).
\end{theorem}

\begin{proof}
Choose the witness \(\left|E(\operatorname{graph}(M))\right|\).
The substitution uses the following argument.
\begin{proofstep}
Fix \(x\) in \(V(M)\).
Assume \(x \in V(M)\).
By \paperlink{premise:53696d706c6547726170682e53756267726170682e636f655f646567726565}{Prerequisite~P\ref{premise:53696d706c6547726170682e53756267726170682e636f655f646567726565}} with \(G^{'} = M\) and \(v = x\), we obtain \(\deg_{\operatorname{graph}(M)}(x) = \deg_{M}(\operatorname{val}(x))\).
By \paperlink{theorem:1}{Lemma~\ref{theorem:1}}, we obtain \(M\text{ is a matching} \iff \left(\forall v\, :\, V,\; v \in V(M) \implies \deg_{M}(v) = 1\right)\).
By propositional extensionality, we obtain \(\left(M\text{ is a matching}\right) = \left(\forall v\, :\, V,\; v \in V(M) \implies \deg_{M}(v) = 1\right)\).
Applying the same function to equal arguments, we obtain \((x_{1} \mapsto x_{1})(M\text{ is a matching}) = (x_{1} \mapsto x_{1})(\forall v\, :\, V,\; v \in V(M) \implies \deg_{M}(v) = 1)\).
Use the hypothesis \(M\text{ is a matching}\).
Rewriting with \paperlink{premise:537562747970652e636f655f70726f70}{Prerequisite~P\ref{premise:537562747970652e636f655f70726f70}} gives \(\operatorname{val}(x) \in V(M)\).
Substituting along the equality, we obtain \(\deg_{M}(\operatorname{val}(x)) = 1\).
Chaining the equalities, we obtain \(\deg_{\operatorname{graph}(M)}(x) = 1\).
\end{proofstep}
Rewriting with \paperlink{premise:5365742e746f46696e7365745f63617264}{Prerequisite~P\ref{premise:5365742e746f46696e7365745f63617264}}, \paperlink{premise:46696e7365742e73756d5f636f6e6772}{Prerequisite~P\ref{premise:46696e7365742e73756d5f636f6e6772}}, \paperlink{premise:46696e7365742e73756d5f636f6e7374}{Prerequisite~P\ref{premise:46696e7365742e73756d5f636f6e7374}}, \paperlink{premise:6d756c5f6f6e65}{Prerequisite~P\ref{premise:6d756c5f6f6e65}} and the identity established above gives \(|\operatorname{as\,finite\,set}(V(M))| = \sum_{v\in V(M)}(\deg_{\operatorname{graph}(M)}(v))\).
Rewriting with \paperlink{premise:65715f6f665f686571}{Prerequisite~P\ref{premise:65715f6f665f686571}}, \paperlink{premise:46696e7365742e73756d5f636f6e6772}{Prerequisite~P\ref{premise:46696e7365742e73756d5f636f6e6772}}, \paperlink{premise:46696e747970652e696e73744661737453756273696e676c65746f6e}{Prerequisite~P\ref{premise:46696e747970652e696e73744661737453756273696e676c65746f6e}} and substitution of equals in the preceding result gives the same displayed equality.
Rewriting with \paperlink{premise:53696d706c6547726170682e73756d5f646567726565735f65715f74776963655f636172645f6564676573}{\emph{the degree-sum formula}} (\paperlink{premise:53696d706c6547726170682e73756d5f646567726565735f65715f74776963655f636172645f6564676573}{Prerequisite~P\ref{premise:53696d706c6547726170682e73756d5f646567726565735f65715f74776963655f636172645f6564676573}}) in the preceding result gives \(\left|\operatorname{as\,finite\,set}(V(M))\right| = 2 \cdot \left|E(\operatorname{graph}(M))\right|\).
Rewriting with \paperlink{premise:74776f5f6d756c}{Prerequisite~P\ref{premise:74776f5f6d756c}} in the preceding result gives \(|\operatorname{as\,finite\,set}(V(M))| = |E(\operatorname{graph}(M))| + |E(\operatorname{graph}(M))|\).
This witness proves existence. Hence \(\operatorname{Even}\left(\left|\operatorname{as\,finite\,set}(V(M))\right|\right)\).
\end{proof}

\end{document}
